\section{Introduction}

\acp{bss} are a crucial part of urban mobility solutions, providing an eco-friendly alternative to private vehicles. Users notably reduce their reliance on other transportation modes~\cite{Ma2020}, leading to increased physical activity and travel time savings~\cite{Qiu2018}. Additionally, urban environments benefit from reduced fuel consumption and improved economic development~\cite{Qiu2018}. Moreover, \acp{bss} have experienced substantial growth since the 2000s, with 2022's global landscape encompassing nearly 2,000 \acp{bss} and an impressive fleet of approximately 9 million bicycles~\cite{MeddinWorldMapReport2022}. The scale of these systems and technological advancements have enabled the generation of large amounts of data, providing researchers and decision-makers with valuable insights to improve existing systems.

\ac{bss} management is often discussed in terms of mobility and rebalancing; however, maintenance is also important for system reliability and user experience. While the importance of maintenance as a critical activity for companies has been underscored in various studies~\cite{Cho1991, Reinertsen1996}, \ac{bss} maintenance practices have predominantly relied on the traditional approach of applying corrective measures (i.e., addressing issues as they arise by replacing worn-out or damaged parts). Therefore, enhancing the \ac{bss} user experience and operational efficiency could be achieved by developing a \ac{pm} strategy capable of predicting bike component failures before they occur, allowing for timely scheduling of maintenance activities. This approach would not only improve the availability and security of the bike fleet but also reduce costs.

To address this challenge, this Chapter leverages both trip and maintenance datasets provided by \textit{Bicing}. Our research revolves around three key questions. First, we explore bike component failure patterns, seeking to understand how the bike model factor contributes to the wear and tear of distinct parts. Second, we aim to assess the potential of forecasting bike component failures by evaluating the accuracy of various survival models. Specifically, our \ac{pm} system focuses on three key components: brake pads, wheel spokes, and chains, chosen for their substantial representation in the data. Third, we aim to identify the critical factors that influence predicting the longevity of the bike components.


\section{Literature review}

Since the 1990s, the scientific literature has emphasized the importance of maintenance as a critical activity for companies to improve reliability and reduce costs~\cite{Cho1991, Reinertsen1996}. Over time, knowledge and techniques have evolved, transitioning from a corrective maintenance, which involved addressing failures after they occurred, to a preventive maintenance~\cite{Ran2019}, which implements scheduled maintenance based on time intervals to prevent breakdowns. While this approach helps to mitigate most failures, it comes with the drawback of high prevention costs. With the increased computational power, the advancement of \ac{ai}, and the rise of the IoT, a new strategy called \ac{pm} has emerged, which can predict failures before they happen, allowing an even lower failure rate and reduced costs. One of its main challenges is fault prognosis, which focuses on forecasting when a failure is likely to occur~\cite{Wang2017_paradigm}. By accurately predicting failures, maintenance activities can be scheduled in advance, minimizing downtime, and optimizing resource allocation.

Reliability theory plays a significant role in \ac{pm} modeling, and it is essentially the same as \ac{sa}~\cite{Yang2022_prognostic}. \ac{sa} was originally developed in the field of biomedical sciences to examine life tables~\cite{Cox1972}, but its concept of events can be applied to various domains, such as machine component failures. A key challenge in \ac{sa} studies is censoring, which refers to missing data when an event is not observed~\cite{Leung1997, Gijbels2010}. Censoring occurs not due to technical failures but rather due to the nature of the studied event. For instance, if a participant in a clinical trial decides to discontinue his participation before the event of interest occurs. It is precisely the presence of censored data that makes impractical the application of predictive algorithms using the usual statistical and \ac{ml} approaches.

According to~\cite{Wang2019_ml_review}, survival methods can be classified into statistical and \ac{ml} methods. Statistical methods focus on characterizing the distribution of event times and the statistical properties of parameter estimation, such as estimating survival curves. These first models can be further divided into: (1) Non-parametric models (Kaplan-Meier, Nelson-Aalen, Life-Table), which make no assumptions about the underlying distribution and estimate the survival curves directly from the data. (2) Semi-parametric models (Cox model, CoxBoost, Time-dependent Cox), which incorporate both non-parametric estimation of the baseline survival function and parametric estimation of the effects of covariates. (3) Parametric models (Penalized regression, Accelerated Failure Time), which assume a specific distribution for the survival time and estimate the parameters of that distribution. \ac{ml} methods combine traditional \ac{sa} techniques with \ac{ml} algorithms, such as survival trees, Bayesian methods, neural networks, or support vector machines. Advanced \ac{ml} techniques, including ensemble learning, active learning, transfer learning, and multitask learning, have also been applied in the field of \ac{sa}.

These methodologies have been applied in a wide range of fields, including healthcare~\cite{Reddy2015}, reliability~\cite{Modarres2016}, crowdfunding~\cite{Li2016}, student retention~\cite{Ameri2016}, customer lifetime~\cite{Furrer2022}, and unemployment duration analysis~\cite{Kiefer1988}. However, to the best of our knowledge, there are no applications of \ac{sa} to \acp{bss}. Apart from \ac{sa}, scientific literature on bike's \ac{pm} is scarce.~\cite{MountainBike2019} proposed using smartphone vibration readings and support vector machine models to predict the health of mountain bike's components (the rotor, the chain, the wheel bearings, the steering head, and the derailleur cog). However, the scope of this study was not a fleet of bikes but a single bicycle. On the contrary,~\cite{Predictivemaintenancebike2018}'s main objective was to study the cyclists' behavioral patterns in the \ac{bss} of Oslo, Norway, and successfully identify the need for bike maintenance. In this case, random forest models were applied to the ride (destination, duration, and date) and the cyclist (gender and year of birth) data. Finally,~\cite{Matkovic2021} focuses on predicting brakes' performance with \ac{knn}, \ac{lstm} and \ac{xgb} classifiers, using as input physical influences and acceleration/deceleration forces coming from hall and inertial \ac{iot} sensors.